\chapter{Vorgaben und gleichzeitige Auswahlen}
\label{b09guide:ch:selections}

Hier geht es ausschließlich um die Existenz gemeinsamer Auswahlen.
Alle Voraussetzungen betreffen Mengen, ihre Elemente und
Teilmengen. Jeder Satz gilt unter AC; seine Aussage für alle
zulässigen Mengendaten ist umgekehrt äquivalent zu AC über ZF.

\section{Fortsetzung einer partiellen Auswahl}

\begin{definition}[Partielle Auswahl]
Für eine Familie $(A_i)_{i\in I}$ und $J\subseteq I$ heißt
\[
 p\in\prod_{j\in J}A_j
\]
eine \emph{partielle Auswahl mit Definitionsbereich $J$}.
Eine \emph{Fortsetzung} ist ein $f\in\prod_{i\in I}A_i$
mit $f\restriction_J=p$, also $f(j)=p(j)$ für jedes $j\in J$.
\end{definition}

\begin{theorem}[Fortsetzungssatz für Auswahlen]
\label{b09guide:thm:extension}
Seien alle $A_i$ nichtleer, $J\subseteq I$ und
$p\in\prod_{j\in J}A_j$ gegeben. Unter AC gibt es
\[
 f\in\prod_{i\in I}A_i
 \quad\text{mit}\quad f\restriction_J=p.
\]
Jede bereits festgelegte partielle Auswahl läßt sich also
auf die ganze Indexmenge ausdehnen, ohne einen ihrer Werte zu ändern.
Die allgemeine Aussage ist äquivalent zu AC.
\end{theorem}

Dies schließt $J=\varnothing$ und $J=I$ ein. Eine Einschränkung auf
endlich viele vorgegebene Werte ist nicht erforderlich.

\section{Auswahl mit zulässigen Werten}

\begin{theorem}[Mengenvorgaben und vorgeschriebene Werte]
\label{b09guide:thm:constraints}
Seien $(A_i)_{i\in I}$ und $(B_i)_{i\in I}$ Familien mit
$B_i\subseteq A_i$, ferner $J\subseteq I$ und
$p\in\prod_{j\in J}A_j$.
Unter AC sind folgende Bedingungen äquivalent:
\begin{enumerate}
\item Es gibt $f\in\prod_{i\in I}A_i$ mit
      $f\restriction_J=p$ und $f(i)\in B_i$ für alle $i\in I$.
\item Für alle $j\in J$ gilt $p(j)\in B_j$, und für alle
      $i\in I\setminus J$ gilt $B_i\ne\varnothing$.
\end{enumerate}
Die Aussage für alle solchen Daten ist äquivalent zu AC.
\end{theorem}

Die Mengen $B_i$ beschreiben genau die zugelassenen Werte.
Sie dürfen etwa durch zusätzliche Zugehörigkeitsbedingungen,
Schnitte oder das Ausschließen verbotener Elemente festgelegt sein.
Vorausgesetzt wird die Nichtleerheit der \emph{gesamten} zulässigen
Menge an jedem noch freien Index. Paarweise verträgliche
Einzelbedingungen allein ersetzen diese Voraussetzung nicht.

\section{Gleichzeitige Auswahlfunktionen}

\begin{theorem}[Simultane Auswahl auf einer Familie von Grundmengen]
\label{b09guide:thm:simultaneous}
Für jede indizierte Familie von Mengen $(U_t)_{t\in T}$ gibt es
unter AC eine Familie $(c_t)_{t\in T}$ mit
\[
 c_t\in\ChoiceSet(U_t)\qquad(t\in T).
\]

Auch diese allgemeine Aussage ist äquivalent zu AC.
\end{theorem}

Die Indexmenge $T$ und jedes $U_t$ dürfen leer sein.
Die Auswahlfunktionen bilden eine gemeinsame Familie.
Eine Übereinstimmung auf gemeinsamen Teilmengen wird nicht verlangt.

\clearpage
\section{Disjunkte Verfeinerung einer Überdeckung}

\begin{definition}[Indizierte Überdeckung und disjunkte Verfeinerung]
Eine Familie $(A_i)_{i\in I}$ von Teilmengen von $U$ heißt
\emph{Überdeckung} von $U$, wenn $\bigcup_{i\in I}A_i=U$.
Eine \emph{indizierte disjunkte Verfeinerung} ist eine Familie
$(B_i)_{i\in I}$ mit
\[
 \begin{gathered}
 B_i\subseteq A_i\quad(i\in I),\qquad
 \bigcup_{i\in I}B_i=U,\\
 B_i\cap B_k=\varnothing\quad(i,k\in I,\ i\ne k).
 \end{gathered}
\]
Einige oder alle $B_i$ dürfen leer sein.
Die Zuordnung zum ursprünglichen Index ist Teil der Vorgabe.
\end{definition}

\begin{theorem}[Disjunkte Verfeinerung unter Beibehaltung der Indizes]
\label{b09guide:thm:refinement}
Unter AC besitzt jede indizierte Überdeckung einer Menge
eine indizierte disjunkte Verfeinerung.
Die Aussage für alle Mengen $U$ und alle ihre indizierten
Überdeckungen ist äquivalent zu AC.
\end{theorem}

Jedes Element von $U$ gehört danach zu genau einem der Stücke
$B_i$, das innerhalb des zugehörigen $A_i$ liegt.
Die Nichtleerheit jedes Stücks wird ausdrücklich nicht verlangt.

